% !TeX root = ../main.tex
\section{What is implemented}
\label{sec:impl}

All code is in \file{guitar_core/src/tone/}. It is blocking and CPU-heavy, so it
runs on worker threads only. This section follows the data through it, with the
constants the code uses, so that a change in the code can be checked against
this description.

\subsection{Overview}

\begin{inapp}[: entry points]
\begin{itemize}
\item \code{tone::estimate(target, dry, &Options)} \done: \code{target} and
  \code{dry} are 48\kHz{} mono and aligned. Returns an \code{Estimate}: the best
  search result (\code{best}), the closest factory preset (\code{preset}), the
  template settings, the excerpt used, the measured echo (\wip) and the number
  of renders.
\item \code{tone::estimate_from_notes(target, notes, &Options)} \done: renders the
  stand-in dry signal from notes with the plucked-string model, then calls
  \code{estimate}.
\item \code{tone::estimate_from_transcription(audio, notes, &Options)} \wip: the
  same, from the transcriber's notes (the model's mean posterior stands in for
  velocity: $v = 0.4 + 0.6\,a$, clamped to $[0.05, 1]$).
\item \code{tone::save_preset(estimate, title, artist, track)} \wip: writes the
  match IR and a user preset (Section~\ref{sec:impl-save}).
\end{itemize}
\code{Options} defaults: excerpt 12\,s, population 16, 30 generations, all CPU
threads, seed 1, spectral floor 80\dB.
\end{inapp}

The algorithm, in order:
\begin{enumerate}
\item Cut the 12\,s window of the target with the most energy (energy summed in
  0.1\,s blocks; \code{loudest_window}), and the same window of the dry signal.
\item Analyse the target once (\code{features::Target}).
\item Measure the delay time, if there is an echo (\code{analysis::echo_delay},
  \wip).
\item Score every factory preset at input RMS 0.1 (presets with a Looper or a
  Pitch Shifter are excluded); keep the closest as the baseline.
\item For each amp voicing (Clean, Crunch, Lead), run separable CMA-ES over the
  template chain from a neutral start; keep the best of the three.
\item Return the best chain and the EQ it still needs, which becomes the match
  IR.
\end{enumerate}

\subsection{The stand-in dry signal}
\label{sec:impl-dry}

\begin{keyidea}[: Karplus--Strong]
A plucked string can be simulated with a delay line whose length is one period
of the note, filled with a short noise burst and fed back through a gentle
low-pass filter: each pass around the loop the burst repeats (giving the pitch)
and loses some treble (giving the natural decay)~\cite{karplus1983digital}. The
\emph{extended} version~\cite{jaffe1983extensions} adds a fractional delay for
exact tuning, a decay-time control and a pluck-position comb filter.
\end{keyidea}

\begin{inapp}[: the plucked-string model, \file{tone/pluck.rs} \done]
For each note (start, end, MIDI pitch, velocity):
\begin{itemize}
\item \textbf{Delay line} of $f_s/f_0 - 0.5$ samples (the two-point loss filter
  adds half a sample), with linear interpolation for the fractional part.
\item \textbf{Excitation}: white noise through a one-pole low-pass (coefficient
  \code{brightness}, plus \code{velocity_brightness} times the velocity), minus a
  copy delayed by \code{pluck_position} times the period (the comb notches a pick
  at that point along the string creates), scaled by velocity$^{1.5}$.
\item \textbf{Loop}: two-point average (high-frequency loss) times a gain $g$
  chosen so the note decays by 60\dB{} in $T_{60}$:
  $g = 10^{-3/(f_0 T_{60})}$. While the note is held,
  $T_{60} = T_{\mathrm{lowE}}\cdot(f_{E2}/f_0)^{0.7}$ with $T_{\mathrm{lowE}}$ = \code{low_e_t60}, clamped to
  $[0.4\,\text{s}, T_{\mathrm{lowE}}]$ (higher notes ring shorter); after the note
  ends, $T_{60}$ = \code{release_secs}.
\item \textbf{Post}: a DC blocker, then the fixed filter \code{DI_EQ}: the
  measured long-term spectrum of real DI minus this model's, as a minimum-phase
  IR.
\end{itemize}
Settings were fitted by CMA-ES (\code{eval::fit_pluck}) on 14 Guitar-TECHS DI
recordings, minimising the EQ-forgiven distance to the real DI of the same
notes: \code{pluck_position} 0.48, \code{release_secs} 0.45, \code{low_e_t60} 16\,s,
\code{brightness} 0.008, \code{velocity_brightness} 0. The distance fell from
13.05 to 9.90\dB{} on the fitting set \measured. Refit only on Guitar-TECHS,
never on EGDB (the benchmark).
\end{inapp}

\begin{inapp}[: the sampler, \file{tone/sampler.rs} \wip]
An alternative stand-in built from real recordings: Guitar-TECHS
``allsinglenotes'' DI (every pitch of a guitar, plucked and held; CC BY 4.0).
Each transcribed note gets the recorded note of the nearest pitch, re-pitched by
resampling if needed, scaled by the velocity ratio (clamped to 0.2--4), started
10\ms{} early so the attack is whole, and faded out over 60\ms{} after the note's
end. It has a real pickup's attack and decay, from another guitar. Selected in
the benchmark with \code{GVA_TONE_SAMPLES}; not yet measured.
\end{inapp}

\subsection{The distance}
\label{sec:app-distance}

\file{tone/features.rs} \done. Both signals are first scaled to RMS 1.

\paragraph{Log-band spectrogram.} At each of two resolutions $r$ the signal is cut
into Hann-windowed frames, transformed, and the power is summed in triangular
bands spaced evenly on a log-frequency axis up to 16\kHz:
\begin{center}
\small
\begin{tabular}{@{}lccc@{}}
\toprule
Resolution & FFT size / hop (at 48\kHz) & Lowest band edge & Bands \\
\midrule
coarse time, fine frequency & 2048 / 512 (43\ms{} / 11\ms) & 60\Hz & 40 \\
fine time, coarse frequency & 512 / 128 (11\ms{} / 2.7\ms) & 250\Hz & 20 \\
\bottomrule
\end{tabular}
\end{center}
Band energies in dB are floored at 80\dB{} below the loudest band (the search
may use a higher floor, \code{Options::floor_db}, \wip). Call the target's value
$T_r(t,b)$ and the render's $R_r(t,b)$ for frame $t$ and band $b$.

\paragraph{Active frames.} Only frames of the target within 45\dB{} of its
loudest frame count (set $\mathcal{A}_r$); silence is ignored.

\paragraph{EQ allowance.} The long-term spectra are the means over active frames,
$\bar T_r(b)$ and $\bar R_r(b)$. Their difference
$\Delta_r(b) = \bar T_r(b) - \bar R_r(b)$ is split into its mean
$o_r = \overline{\Delta_r}$ (a level offset) and a centred, capped EQ
\[
  e_r(b) = \operatorname{clip}\bigl(\Delta_r(b) - o_r,\ -18, +18\bigr)\ \text{dB}.
\]
\paragraph{Search distance.}
\[
  D = \frac{1}{2}\sum_{r=1}^{2} \frac{1}{|\mathcal{A}_r|\,B_r}
      \sum_{t\in\mathcal{A}_r}\sum_{b=1}^{B_r}
      \bigl|\,T_r(t,b) - \bigl(R_r(t,b) + e_r(b) + o_r\bigr)\bigr|.
\]
In words: the average number of decibels by which the render misses the target
in each band of each active frame, \emph{after} the render has been given the
best smooth EQ (within $\pm18\dB$) that matches the long-term spectra.

\paragraph{Benchmark score} (\code{plain_distance}): the same formula without
$e_r$ and $o_r$, and always with the 80\dB{} floor. It measures how far apart
the two signals sound, colour included.

\begin{keyidea}[: why forgive an EQ during the search]
The stand-in dry signal has the right notes but the wrong guitar, and the target
went through a cabinet, a microphone and a mix EQ that the template chain does
not model. Both are mostly \emph{linear} differences, which a filter fixes
afterwards (Section~\ref{sec:dsp}). Forgiving them lets the search spend its
effort on what only it can fix: input level, compression, drive, amp gain and
time-based effects. Distortion is \emph{not} forgiven: clipping changes the
spectrum frame by frame, not by a fixed EQ (a unit test checks this:
\code{distortion_is_not_forgiven}).
\end{keyidea}

\subsection{The template chain}

\file{tone/template.rs}. The search moves a vector $x\in[0,1]^n$; each
coordinate maps to a pedal parameter, linearly or on a log scale
(Table~\ref{tab:template}). Order: compressor $\to$ distortion $\to$ amp $\to$
delay $\to$ reverb. A pedal whose amount (compressor amount, drive mix, delay
mix, reverb mix) is below 0.04 is left out of the chain, so the search can
switch pedals off.

\begin{table}[h]
\centering
\caption{Searched parameters. The first ten are in the committed version
\done; the two delay parameters are in the working tree \wip.}
\label{tab:template}
\small
\begin{tabular}{@{}rllll@{}}
\toprule
\# & Parameter & Range & Scale & Maps to \\
\midrule
0 & input level (RMS) & 0.01--0.4 & log & level the dry signal is played at \\
1 & compressor amount & 0--1 & lin & ratio $1 + 7\cdot$amount (up to 8:1) \\
2 & compressor threshold & 0.02--0.5 & log & threshold \\
3 & drive mix & 0--1 & lin & Distortion mix \\
4 & drive & 0.5--10 & log & Distortion drive \\
5 & drive tone & 0.15--1 & lin & Distortion tone \\
6 & amp gain & 0.5--5 & log & Amp gain (voicing searched separately) \\
7 & reverb mix & 0--0.5 & lin & Reverb mix \\
8 & reverb room & 0--0.95 & lin & Reverb room size \\
9 & reverb damping & 0--0.95 & lin & Reverb damping \\
10 & delay mix & 0--0.5 & lin & Delay mix (time measured, not searched) \\
11 & delay feedback & 0--0.7 & lin & Delay feedback \\
\bottomrule
\end{tabular}
\end{table}

Fixed values: compressor attack 10\ms, release 150\ms, make-up gain 1;
distortion level 1, bias 0; amp tone stack flat (bass, mid, treble 0), because
it acts after the clipping and is therefore linear and left to the match IR;
delay with one tap and damping 0.3. The neutral start is input RMS 0.1,
compressor amount 0.2, drive mix 0.3, reverb mix 0.15, delay mix and feedback 0.3,
everything else 0.5.

\subsection{The search}

\file{tone/cmaes.rs} implements separable CMA-ES~\cite{ros2008sepcmaes}
(Section~\ref{sec:search}) with a deterministic random generator (SplitMix64), so
the same seed gives the same result. Candidates outside the unit cube are
evaluated at the nearest point inside, plus a penalty of 10 times the squared
distance to the cube. Each generation's candidates are rendered in parallel.
Settings: population 16, 30 generations, initial step size 0.3, seed
incremented per voicing. The amp voicing is discrete, so there is one search per
voicing.

\subsection{Measuring the delay time}
\label{sec:impl-delay}

\file{tone/analysis.rs} \wip.
\begin{keyidea}[: echoes repeat onsets]
A delay repeats every onset of the target after the delay time. The stand-in dry
signal has the same notes without echoes. So a time lag at which the target's
onset pattern correlates with itself much more than the dry signal's does is the
delay time; rhythmic repetition in the playing shows up in both and cancels.
\end{keyidea}
Both signals are turned into an onset envelope (positive spectral flux of the log
magnitude, 1024-point FFT, one value every 5\ms). For each lag from 60\ms{} to
1.2\,s the code computes the normalised autocorrelation of the target envelope
minus that of the dry envelope (the \emph{excess}). If the largest excess is
below 0.12 there is no echo and no delay pedal. Otherwise it takes the
\emph{shortest} lag whose excess reaches 70\% of the best (a feedback delay also
correlates at two and three times its time), refined to the local maximum. The
benchmark counts a delay time within 5\% of the truth as right.

\subsection{The match filter}
\label{sec:impl-match}

\file{tone/matching.rs} \done. The EQ $e(b)$ of the finer-frequency resolution
(40 bands) is turned into an impulse response:
\begin{enumerate}
\item Interpolate the EQ on a log-frequency axis at every bin of an 8192-point
  FFT (flat beyond the outer bands).
\item Take the natural log of the magnitude and its inverse FFT: the real
  cepstrum.
\item Fold it: keep coefficient 0 and $N/2$, double the causal half, zero the
  rest. This is the cepstrum of the minimum-phase filter with the same
  magnitude~\cite{oppenheim2010dtsp}.
\item FFT, complex exponential, inverse FFT: the minimum-phase IR. Keep the
  first 2048 samples (43\ms) and fade the last quarter with a half-cosine so the
  cut does not ring.
\end{enumerate}
The IR is written as a WAV file and loaded by a Cabinet pedal (mix 1) appended at
the end of the chain (\code{Candidate::chain_with_ir}).

\subsection{Saving the result and the UI}
\label{sec:impl-save}

\wip{} (working tree, not committed at the time of writing).
\code{save_preset} writes the IR to \file{<data dir>/guitar_virtual_amp/tone_match/}
and a user preset named ``\emph{title} -- \emph{track} (timbre estimado)'' with
\code{song} and \code{artist} set and the tags \code{timbre-estimado} and
\code{faixa:<track>}; estimating again overwrites it.
\code{shared::tab::tone::preset_fits_track} keeps a preset tagged for one track
from being suggested for another. In the Transcribe form, ``Estimar o timbre
(experimental)'' (\code{AppSettings::transcription_tone}, off by default, needs
separation) adds a last stage to \code{transcribe::run} that estimates each
guitar part from its own audio (its pan region of the stem) and notes; the form
says it takes about 30\,s per guitar.
